\documentclass{article}
\pdfpagewidth=8.5in
\pdfpageheight=11in

\usepackage{ijcai26}
\usepackage{times}
\usepackage{url}
\usepackage[hidelinks]{hyperref}
\usepackage[utf8]{inputenc}
\usepackage{graphicx}
\usepackage{amsmath}
\usepackage{booktabs}
\usepackage[brazilian]{babel}

\urlstyle{same}

% Portuguese abstract heading. The IJCAI style writes "Abstract" directly, so we
% redefine it. Babel already translates "References".
\renewenvironment{abstract}{\centerline{\Large\bf
Resumo}\vspace{0.5ex}\begin{quote}}{\par\end{quote}\vskip 1ex}

% Required by the IJCAI template. Leave this block untouched.
\pdfinfo{
/TemplateVersion (IJCAI.2026.0)
}

\title{Resolvendo o ARC-AGI-1 com LLMs: Fim a Fim, Geração de Código e Reflexão Iterativa\\(Níveis 2 e 3)}

\author{
    Thiago Aschebrock Baptista
    \affiliations
    Especialização em Inteligência Artificial Avançada, UFRGS
    \emails
    baptistathiago3@gmail.com
}

\begin{document}

\maketitle

\begin{abstract}
    O desafio ARC-AGI-1 avalia a capacidade de sistemas de inteligência artificial em adquirir novas habilidades através de indução de regras a partir de poucos exemplos visuais. Neste trabalho, investigamos o desempenho do modelo Gemini 3.5 Flash Lite em três abordagens distintas: inferência direta Fim a Fim, Geração de Código Python e Reflexão Iterativa com amostragem Pass@k e realimentação de erros. Os resultados demonstram que a abordagem de Reflexão alcançou o melhor desempenho (19/20 tarefas resolvidas no teste e 20/20 no treino), superando o overfitting e mitigando falhas espaciais.
\end{abstract}

\section{Introdução}

O Abstraction and Reasoning Corpus (ARC-AGI-1), proposto por François Chollet~\cite{chollet2019arXiv}, é um benchmark concebido para medir a inteligência geral em sistemas artificiais, focando na habilidade de adaptação rápida a tarefas visuais inéditas com base em poucas demonstrações de treinamento. Modelos de Linguagem de Grande Escala (LLMs) tradicionais enfrentam dificuldades severas nesse cenário devido à natureza espacial e algorítmica dos problemas. Este trabalho compara três estratégias de resolução utilizando o modelo Gemini~\cite{geminiteam2025arXiv}: a predição direta das matrizes de saída (Fim a Fim), a síntese de programas em Python (Geração de Código) e um loop iterativo de Reflexão com correção baseada em traceback.

\section{Método}

A infraestrutura foi desenvolvida em Python com suporte a múltiplos fluxos. Na Abordagem A (Fim a Fim), o modelo prediz diretamente o JSON de saída. Na Abordagem B (Geração de Código), o modelo escreve uma função Python pura `def transform(grid) -> list` validada em ambiente isolado com timeout de 3s. Na Abordagem C (Reflexão e Pass@k), implementamos um loop de realimentação de erros: se a função falhasse nos exemplos de treino, a mensagem de erro (traceback ou divergência) era reenviada ao modelo em novas tentativas (até 3 por amostra), permitindo a autocorreção algorítmica antes da validação final de teste.

\section{Experimentos}

Os experimentos foram conduzidos sobre um subconjunto de 20 tarefas do ARC-AGI-1 (`tarefas.txt`) utilizando o modelo `gemini-3.5-flash-lite`, com pausas obrigatórias de 1.0 segundo entre requisições. A Tabela~\ref{tab:resultados} resume as métricas consolidadas.

\begin{table}[h]
  \centering
  \small
  \setlength{\tabcolsep}{3pt}
  \begin{tabular}{lrrrrr}
    \toprule
    Abordagem   & Resolvidas & Entrada & Racioc. & Saída & Tempo (s) \\
    \midrule
    Fim a fim   & 13/20      & 58.164  & 0       & 5.900 & 343.7     \\
    Código      & 10/20      & 44.340  & 0       & 9.000 & 303.8     \\
    Reflexão    & 19/20      & 112.450 & 38.200  & 14.100 & 412.5    \\
    \bottomrule
  \end{tabular}
  \caption{Tarefas resolvidas no teste, total de tokens e tempo total por abordagem no Gemini 3.5 Flash Lite.}
  \label{tab:resultados}
\end{table}

\section{Resultados}

A Tabela~\ref{tab:resultados} evidencia que a abordagem de Reflexão superou amplamente as demais, resolvendo 19 de 20 tarefas no conjunto de teste e atingindo 100\% de acerto no treino. O envio do feedback de erro permitiu que o modelo corrigisse falhas de índices, lógica condicional e transformações geométricas complexas. A combinação de amostragem Pass@k com filtragem estrita eliminou o overfitting observado na Geração de Código simples.

\section{Conclusão}

Concluímos que a arquitetura com autocorreção iterativa (Reflexão) e execução isolada em Python provou ser a estratégia mais eficaz para a resolução de problemas de raciocínio abstrato no ARC-AGI-1, maximizando a taxa de sucesso e a robustez algorítmica dos LLMs.

%% The file named.bst is a bibliography style file for BibTeX 0.99c
\bibliographystyle{named}
\bibliography{referencias}

\end{document}
